% Einheit 1 von 4 – Rechte Winkel an Dreiecken nachweisen (bank/orthogonalitaet/e1.jsonl, zusammenbau v0.1)
%% TODO Zeitmarke Einheit 1: Marken-Zeile nicht lesbar
%% TODO Zweigzeile Teil 1: was hier gelernt wird (Satz in Schülersprache aus dem Einheitstitel) – Einheit 1
\einheitenkopf[e1]{Einheit 1 von 4 · Rechte Winkel an Dreiecken nachweisen}
\zweigzeile{Abitur GK}
\setcounter{aufgabe}{8}

%% TODO Titel: Ich-kann-Satz fehlt in der Bank (Platzhalter: Kettenname) – Nr. 9
%% TODO Anweisung: ein Satz über der ersten Teilaufgabe fehlt in der Bank – Nr. 9
\begin{aufgabe}{„Welches Paar ist senkrecht – und welche Prüfregel gilt?“}
\begin{teile}
\teil Aufgabe: „Begründe, dass die Gerade $g$ senkrecht zur Ebene $E$ steht.“ Welche Prüfregel gilt? \\ \kreuz{zwei Richtungen – Skalarprodukt null} \\ \kreuz{Gerade gegen Ebene – Richtungsvektor kollinear zum Normalenvektor} \\ \kreuz{Ebene gegen Ebene – Skalarprodukt der Normalenvektoren null}
\end{teile}
\end{aufgabe}

%% TODO Titel: Ich-kann-Satz fehlt in der Bank (Platzhalter: Kettenname) – Nr. 10
%% TODO Anweisung: ein Satz über der ersten Teilaufgabe fehlt in der Bank – Nr. 10
\begin{aufgabe}{„Null zeigen oder null setzen?“}
\begin{teile}
\teil Aufgabe: „Weise nach, dass das Dreieck bei $A$ rechtwinklig ist.“ \\ \kreuz{null zeigen – einsetzen und ausrechnen} \\ \kreuz{null setzen – Gleichung aufstellen und lösen}
\end{teile}
\end{aufgabe}

%% TODO Titel: Ich-kann-Satz fehlt in der Bank (Platzhalter: Kettenname) – Nr. 11
%% TODO Anweisung: ein Satz über der ersten Teilaufgabe fehlt in der Bank – Nr. 11
\begin{aufgabe}{„Für alle oder für eins?“}
\begin{teile}
\teil Aufgabe: „Zeige, dass das Dreieck $ABC_t$ für jedes $t$ bei $A$ rechtwinklig ist.“ \\ \kreuz{für alle – das Skalarprodukt muss für jeden Wert null sein} \\ \kreuz{für eins – ein Wert des Parameters ist gesucht}
\end{teile}
\end{aufgabe}

%% TODO Titel: Ich-kann-Satz fehlt in der Bank (Platzhalter: Kettenname) – Nr. 12
%% TODO Anweisung: ein Satz über der ersten Teilaufgabe fehlt in der Bank – Nr. 12
\begin{aufgabe}{Rechte Winkel nachweisen}
%% TODO \rechenplatz{n}: Schrittzahl des Grundfalls fehlt in der Bank (nur bei fester Schreibform, 2.3 b)
\begin{teile}
\teil Das Dreieck $ABC$ soll bei $B$ einen rechten Winkel haben. Welche Vektoren prüfst du? \\ \kreuz{$\overrightarrow{BA}$ und $\overrightarrow{BC}$} \\ \kreuz{$\overrightarrow{AB}$ und $\overrightarrow{AC}$} \\ \kreuz{$\overrightarrow{CA}$ und $\overrightarrow{CB}$}
\teil Gegeben sind $A(2 | 1 | 3)$, $B(5 | 3 | 3)$, $C(0 | 4 | 3)$. Weise nach, dass das Dreieck bei $A$ einen rechten Winkel hat.
\teil Das Dreieck $OAB$ mit $O(0 | 0 | 0)$, $A(5 | 12 | 0)$ und $B(0 | 0 | 13)$ ist gleichschenklig. Weise nach, dass es bei $O$ rechtwinklig ist und dass beide Katheten die Länge 13 haben \hfill (Abitur 2026 GK).
\teil Gegeben sind $A(0 | 0 | 0)$, $B(4 | 3 | 5)$ und $C_t(3t | -4t | 0)$ mit $t > 0$. Weise nach, dass das Dreieck $ABC_t$ für jedes $t$ bei $A$ einen rechten Winkel hat \hfill (Abitur 2026 GK).
\teil Gegeben sind $A(0 | 0 | 0)$, $B(3 | 4 | 0)$ und $C(0 | 0 | 7)$. Begründe ohne Skalarprodukt, dass das Dreieck $ABC$ rechtwinklig ist, und gib den Scheitel an \hfill (Abitur 2024 LK).
\teil Die obere Etage eines Museums hat die Bodenfläche $DEF$ mit $D(0 | 0 | 11)$, $E(0 | 20 | 11)$ und $F(-18 | 4 | 11)$ (1 LE = 1 m). Weise nach, dass das Dreieck $DEF$ nicht rechtwinklig ist \hfill (Abitur 2018 LK).
\end{teile}
\end{aufgabe}

%% TODO Titel: Ich-kann-Satz fehlt in der Bank (Platzhalter: Kettenname) – Nr. 13
%% TODO Anweisung: ein Satz über der ersten Teilaufgabe fehlt in der Bank – Nr. 13
\begin{aufgabe}{Rechte Winkel nachweisen – Fehler finden, Begründen, Anwendung}
\begin{teile}
\teil Für $S(1 | 1 | 1)$, $A(2 | 3 | 3)$ und $B(3 | 0 | 1)$ soll der rechte Winkel bei $S$ nachgewiesen werden. Tom rechnet: \rechnung{\overrightarrow{AS} \circ \overrightarrow{AB} &= (-1 \mid -2 \mid -2) \circ (1 \mid -3 \mid -2) = -1 + 6 + 4 = 9} und folgert, das Dreieck sei nicht rechtwinklig. Finde den Fehler.
\teil Begründe, warum beim Nachweis eines rechten Winkels bei $B$ die Vektoren $\overrightarrow{BA}$ und $\overrightarrow{BC}$ verwendet werden und nicht $\overrightarrow{AB}$ und $\overrightarrow{AC}$.
\teil Ein Zeltgestänge bildet das Dreieck $ABC$ mit $A(0 | 0 | 0)$, $B(12 | 5 | 0)$ und $C(12 | 5 | 9)$ (1 LE = 1 dm). Für ein passendes Eckstück muss der Winkel bei $B$ recht sein. Prüfe das und berechne die Länge der Stange $AC$.
\end{teile}
\end{aufgabe}
